\chapter{Các pipeline dữ liệu}

Hệ thống có bốn đường ống dữ liệu độc lập. Nhầm lẫn giữa chúng là nguồn gốc của nhiều lỗi,
nên chương này bắt đầu bằng việc phân biệt rõ vai trò.

\begin{longtable}{L{3.4cm} L{4.6cm} L{5.6cm}}
\toprule
\textbf{Pipeline} & \textbf{Đầu vào} & \textbf{Đầu ra}\\
\midrule
\endhead
Content ETL & Kho ngữ liệu mở (OEWN, CMUdict\ldots) & Bản ghi \textbf{từ vựng và phát âm} — \textbf{không bao giờ} là bài học\\
Content Agent & Danh sách từ có nhãn CEFR & Khoá học $\rightarrow$ chương $\rightarrow$ bài học kèm bài tập, ở trạng thái bản nháp\\
Learner Observation & Lượt chat, bài làm & Trạng thái khái niệm và điểm kỹ năng\\
Content Prefetch & RSS, YouTube, podcast & Nội dung thực tế đã lưu đệm\\
\bottomrule
\caption{Bốn pipeline dữ liệu}
\end{longtable}

\section{Content ETL — ngữ liệu từ vựng có kiểm soát bản quyền}

\subsection{Nguồn và giấy phép}

Thư mục \code{ai-service/api/services/content\_etl/} triển khai ETL cho bảy nguồn, mỗi
nguồn có một bộ tiếp hợp (\en{adapter}) riêng:

\begin{longtable}{L{2.8cm} L{3.4cm} L{3.0cm} L{4.0cm}}
\toprule
\textbf{Nguồn} & \textbf{Nội dung} & \textbf{Giấy phép} & \textbf{Trạng thái}\\
\midrule
\endhead
OEWN & Định nghĩa, quan hệ nghĩa & CC BY 4.0 & Bật, ghim bản 2025, 185.129 bản ghi\\
CMUdict & Phiên âm & LicenseRef-CMUdict & Bật, ghim \code{74790861\ldots}, 126.052 bản ghi\\
CEFR-J & Nhãn bậc CEFR & Thương mại & \textbf{Tắt có chủ đích}\\
Wikidata & Thực thể, chủ đề & Mở & Cần ghim mới bật\\
Tatoeba & Câu ví dụ & Mở & Cần ghim mới bật\\
LibriSpeech & Âm thanh đọc sách & Mở & Cần ghim mới bật\\
Common Voice & Âm thanh cộng đồng & Mở & Cần ghim mới bật\\
\code{admin\_upload} & Danh sách từ do quản trị viên tải lên & Thuộc sở hữu & Bật\\
\bottomrule
\caption{Nguồn ngữ liệu và trạng thái ghim}
\end{longtable}

\subsection{Nguyên tắc "ghim mới chạy"}

Một nguồn chỉ được bật khi \textbf{cả} tham chiếu/phiên bản \textbf{và} tổng kiểm tra
sha256 đều được ghim trong biến môi trường. Nguồn chưa ghim bị bỏ qua trong im lặng chứ
không chặn khởi động dịch vụ. Lệnh chạy:

\begin{lstlisting}[language=bash]
python -m api.services.content_etl.cli sync \
    --sources cmudict,oewn --write --storage-root /duong/dan/ngoai/container
\end{lstlisting}

\begin{luuy}[Dữ liệu mẫu dễ dãi hơn dữ liệu thật]
Cả ba bộ tiếp hợp đầu tiên đều xuất xưởng ở trạng thái hỏng và cùng báo "0 bản ghi được
duyệt", vì tệp mẫu trong kiểm thử không giống tệp tải về thật: \code{cmudict.dict} thật
viết chữ thường và có chú thích sau dấu \code{\#} (mẫu thì viết hoa); bản phát hành OEWN
thật được nén gzip và không có \en{namespace} XML (mẫu thì là XML thuần có namespace); và
từ đa âm tiết của OEWN cần mã hoá phần trăm cho \code{source\_url}, nếu không bộ kiểm tra
URL sẽ viết lại giá trị và tổng kiểm tra bản ghi không còn khớp. Bài học: tệp mẫu phải có
\textbf{cùng hình dạng} với tệp thật.
\end{luuy}

\subsection{Vì sao ETL không tự gieo được khoá học}

Đây là hiểu nhầm phổ biến nhất về hệ thống. Hàm \code{plan\_curriculum} chọn từ
\textbf{nghiêm ngặt theo} \code{declared\_cefr}, trong khi bản ghi OEWN và CMUdict
\textbf{không mang nhãn CEFR} — chúng là dữ liệu làm giàu (định nghĩa, ví dụ, phát âm),
không phải danh sách hạt giống. Nguồn duy nhất có nhãn trong sổ đăng ký là CEFR-J, mà
giấy phép của nó là thương mại. Ngoài ra, đính kèm cả kho ngữ liệu sẽ vượt
\code{CONTENT\_AGENT\_MAX\_RECORDS} (20.000, trong khi OEWN có 185.129 bản ghi).

Vì vậy nguồn hạt giống đang dùng thực tế là \code{admin\_upload}: một danh sách từ mang
\code{declared\_cefr} cho từng từ.

\section{Content Agent — sinh khoá học}

\subsection{Kiến trúc hai nửa}

\begin{longtable}{L{5.2cm} L{8.4cm}}
\toprule
\textbf{Phía AI service} & \textbf{Phía backend}\\
\midrule
\endhead
\code{content\_agent/planner.py} — lập chương trình theo bậc CEFR từ danh sách từ. &
\code{content\_agent\_apply.py} — \textbf{đường duy nhất} ghi khoá học, chương, bài học và bài tập vào cơ sở dữ liệu.\\
\code{generator.py} — sinh nội dung bài học và bài tập bằng LLM. &
\code{content\_agent\_validation.py} — làm sạch, kiểm tra hình dạng bài tập.\\
\code{policies.py} — ràng buộc dạng bài tập được phép. &
\code{content\_agent\_jobs.py} — vòng đời tác vụ, trạng thái, thử lại.\\
\code{adapters.py}, \code{store.py} — nối nguồn dữ liệu, lưu tạo tác. &
\code{content\_agent\_uploads.py} — nhận và kiểm tra tệp tải lên.\\
\bottomrule
\caption{Phân chia trách nhiệm của tác tử nội dung}
\end{longtable}

Khoá học sinh ra luôn ở trạng thái \textbf{chưa xuất bản}. Việc xuất bản là quyết định của
con người và bị chặn nếu còn bài học không có bài tập.

\subsection{Mô hình nào thực sự sinh nội dung}

Tác tử chọn theo thứ tự: \code{GroqMissionGenerator} nếu có khoá Groq $\rightarrow$ Gemini
$\rightarrow$ khuôn mẫu tất định. \textbf{Cả hai đường LLM đều rơi về khuôn mẫu cho từng
bài học khi gặp bất kỳ lỗi nào}, nên một khoá API chết trông giống hệt như "sinh nội dung
thành công" — chỉ khác ở chỗ nội dung là câu mẫu kiểu \textit{"Listen, then repeat the
target phrase clearly."} Phải đọc chữ trước khi tin một tác vụ đã chạy đúng.

\subsection{Sinh lại nội dung yếu}

\code{scripts/generate\_exercises\_ai.py --regenerate} viết lại những bài học tuy chơi
được nhưng yếu: dưới ba loại giao diện bài tập khác nhau, không có bài tập sản sinh
(\en{production task}), hoặc dùng câu lệnh chung chung kiểu "Match the words with their
meanings". Nội dung mới chỉ ghi đè sau khi vượt \code{sanitize\_exercises}, nên một lời gọi
thất bại để lại nội dung cũ nguyên vẹn.

\begin{luuy}[Hạn mức tần suất của tầng miễn phí rất chặt]
Cấu hình \code{--concurrency=2 --delay=6} hoàn thành 93/93 bài. Cấu hình
\code{--concurrency=4 --delay=1} mất 93 trong 249 bài vì lỗi 429.
\end{luuy}

\section{Learner Observation — từ luyện tập thành dữ liệu}

Đây là pipeline quan trọng nhất về mặt kiến trúc, vì nó bắc cầu giữa hai dịch vụ mà không
được phép mất dữ liệu.

\subsection{Bốn chặng}

\begin{enumerate}[leftmargin=1.4em]
  \item \textbf{Ghi vùng đệm.} Hội thoại AI ghi hàng loạt (\en{bulk upsert}) các sự kiện
  tất định vào bộ sưu tập MongoDB \code{learner\_observation\_spool} \textbf{trước khi}
  trả về phản hồi không-luồng, hoặc trước dấu kết thúc SSE với phản hồi luồng.
  \item \textbf{Chuyển tiếp có thuê.} Một tiến trình thuê (\en{lease worker}) lấy lô sự
  kiện và gửi sang điểm cuối nội bộ của backend.
  \item \textbf{Ghi outbox rồi mới xác nhận.} Backend ghi dòng \en{outbox} bất biến-khi-lặp
  vào PostgreSQL, xác nhận với tiến trình gửi, rồi mới xử lý.
  \item \textbf{Áp dụng song song.} Nhiều tiến trình áp dụng sự kiện bằng
  \codesp{FOR UPDATE SKIP LOCKED}. Trạng thái khái niệm, \code{state\_epoch} và cờ đã-áp-dụng
  đổi trong \textbf{một giao dịch duy nhất}.
\end{enumerate}

\subsection{Chấm điểm kỹ năng từ hội thoại}

Mỗi lượt chat gắn \textbf{một} quan sát neo mang \code{skill}, \code{score} và
\code{difficulty\_level}. Hàm \code{\_record\_turn\_skill\_evidence} ở
\code{routes/learner\_state.py} chuyển nó thành một \code{ExerciseResult}. Tính bất
biến-khi-lặp dựa hoàn toàn vào cơ chế chống trùng \code{event\_id} của vùng đệm: chỉ những
sự kiện \textbf{mới được chèn} mới được chấm điểm, nên không cần thêm bảng phụ nào.

\begin{luuy}[Danh sách khoá được phép nằm ở hai nơi]
Danh sách khoá hợp lệ của tải trọng quan sát được khai báo cả ở
\code{app/schemas/learner\_state.py} lẫn
\code{ai-service/api/services/learner\_observation\_spool.py}. Một khoá lạ khiến
\textbf{cả lô} bị từ chối. Khi thêm trường mới, phải sửa cả hai.
\end{luuy}

\subsection{Giới hạn an toàn}

Các giới hạn hiện hành là \textbf{ranh giới an toàn}, không phải năng lực đã đo:
60 khái niệm ứng viên mỗi yêu cầu, 100 quan sát mỗi lô gửi backend, hạn chót đọc trạng
thái người học 40 ms, lô dọn dẹp 5.000 dòng. Số liệu p95/p99 thật, mức đồng thời an toàn
và tỉ lệ trúng đệm phải được ghi vào
\code{docs/operations/learner-state-rollout.md} sau các đợt thử tải theo giai đoạn; việc
nâng cấp lên sản xuất bị cấm khi các ô đó còn trống.

\subsection{Phía client}

Các bề mặt có chấm điểm ở Flutter đi qua \code{SkillEventRecorder}
(\code{lib/core/services/skill\_event\_recorder.dart}) theo kiểu bắn-và-quên
(\en{fire-and-forget}) để không chặn giao diện. Hằng số \code{vocabConceptId} ở đây phải
luôn khớp với \code{\_vocab\_concept\_id} trong
\code{backend-service/app/crud/vocabulary.py}.

\section{Content Prefetch — nội dung thực tế}

Ba tác vụ Celery nạp trước nội dung ngoài đời thực để người học không phải chờ nguồn bên
thứ ba: tin tức mỗi 6 giờ, YouTube mỗi 12 giờ, podcast mỗi ngày. Nội dung được lưu đệm
kèm câu hỏi/quiz sinh sẵn; riêng phần quiz mới là thứ tạo ra bằng chứng học tập.

\section{Gieo dữ liệu và kiểm toán}

\subsection{Thứ tự chạy}

\begin{lstlisting}[language=bash]
seed_courses_directly.py -> seed_analytics.py -> seed_demo_data.py -> seed_questions.py
\end{lstlisting}

Trạng thái cơ sở dữ liệu kỳ vọng sau khi gieo đầy đủ: 109 người dùng, 1.614 hoạt động hằng
ngày, 98 câu hỏi. Các kịch bản đều bất biến-khi-lặp: mỗi kịch bản xoá tập người dùng của
chính nó trước rồi mới gieo lại.

\begin{luuy}[Chỉ hai khoá học nhỏ tái tạo được từ kho mã]
\code{seed\_courses\_directly.py} chỉ gieo hai khoá học mẫu trong
\code{app/core/sample\_data\_catalog.py}. Mười ba khoá học lớn (IELTS, Business,
Advanced\ldots) \textbf{chỉ tồn tại trong cơ sở dữ liệu} và không tái tạo được từ kho mã —
chúng đến từ tác tử nội dung.
\end{luuy}

\subsection{Kiểm toán khả năng chơi được}

Sau mỗi lần gieo khoá học, phải xác minh trước khi tin vào ứng dụng:

\begin{lstlisting}[language=bash]
venv/bin/python3 scripts/audit_course_content.py                 # bao cao
venv/bin/python3 scripts/generate_exercises_ai.py                # lap bai hoc rong
venv/bin/python3 scripts/audit_course_content.py --fix-counters  # dong bo bo dem
\end{lstlisting}

Kịch bản kiểm toán có ba chế độ sửa: \code{--fix-counters} (đồng bộ
\code{exercise\_count}), \code{--fix-matching} (sửa danh sách lựa chọn của bài ghép cặp),
\code{--fix-unplayable} (xử lý bài học không chơi được).

\section{Sao lưu như một pipeline}

Dữ liệu người dùng nằm ở đúng hai kho, nên sao lưu chỉ có hai đối tượng:

\begin{longtable}{L{5.0cm} L{3.0cm} L{5.6cm}}
\toprule
\textbf{Đơn vị systemd} & \textbf{Lịch} & \textbf{Hành động}\\
\midrule
\endhead
\code{lexilingo-backup.timer} & 02:30 hằng ngày & Kết xuất cả hai kho, sinh bảng kê sha256, xoá bản cũ hơn 14 ngày.\\
\code{lexilingo-backup-verify.timer} & 04:00 Chủ nhật & Khôi phục bản mới nhất vào container dùng một lần và kiểm tra dữ liệu người dùng.\\
\code{lexilingo-backup-alert@.service} & Khi \code{OnFailure} & Đặt cờ vào \code{/var/lib/lexilingo/}, ghi log mức nghiêm trọng, gọi webhook nếu có.\\
\bottomrule
\caption{Lịch sao lưu}
\end{longtable}

Bốn quy tắc tồn tại vì từng bị vi phạm:

\begin{enumerate}[leftmargin=1.4em]
  \item \textbf{Khôi phục kiểm tra trước khi phá.} \code{restore-prod.sh} kiểm tra tính
  toàn vẹn gzip, tổng kiểm tra bảng kê và phần đầu tệp kết xuất, chụp ảnh dữ liệu sản xuất
  hiện tại vào \code{backups/pre-restore/}, rồi mới xoá lược đồ. Nó in ra chính xác lệnh
  quay lui nếu có sự cố.
  \item \textbf{Một bản kết xuất chưa phải bản sao lưu cho tới khi nó được khôi phục
  thành công.} Đợt kiểm chứng hằng tuần là thứ duy nhất khiến các bản kết xuất hằng ngày
  đáng tin.
  \item \textbf{Ngưỡng kích thước phải theo từng kho và mang tính tương đối.} PostgreSQL
  sản xuất khoảng 1,9 MB sau nén, MongoDB khoảng 79 KB. Một ngưỡng tuyệt đối chung đã từng
  loại bỏ một bản kết xuất MongoDB hợp lệ; ngưỡng đúng là "dưới 50\% bản tốt trước đó".
  \item \textbf{\code{OnFailure} phải đặt trong khối \code{[Unit]}.} Đặt trong
  \code{[Service]} thì systemd bỏ qua và chỉ ghi một dòng log — nghĩa là không ai được
  báo động. Kiểm chứng bằng \codesp{systemctl show -p OnFailure --value <unit>}; kết quả
  rỗng nghĩa là đang hỏng.
\end{enumerate}

\begin{luuy}[Sao lưu hiện vẫn nằm cùng đĩa]
Các bản sao lưu nằm trên cùng ổ đĩa với cơ sở dữ liệu. Cần đặt
\code{OFFSITE\_RCLONE\_REMOTE} cho \code{lexilingo-backup.service} để khắc phục; kịch bản
cảnh báo ở mỗi lần chạy cho tới khi việc này hoàn tất.
\end{luuy}
